\documentclass{book}
\usepackage{times}
\usepackage{graphicx}
\usepackage{lmodern}
\usepackage[utf8]{inputenc}
\usepackage[brazilian]{babel}

\usepackage{geometry}
\geometry{
    margin=1in,
}

\title{Descompilando\\ \large Entendendo o código da máquina}
\author{Anthony Carneiro da Silva}

\begin{document}
\maketitle

\clearpage

\tableofcontents

\clearpage

\chapter{Índice}
\clearpage

\section{Sobre o autor}

Meu nome é Anthony Carneiro da Silva, tenho 20 anos e sou de 
Sobral/CE. Desde cedo, sempre gostei de usar o computador para jogar 
videogame e acessar a internet, ainda cedo tive contato com o mundo 
da programação, mas assim como muitos na época acreditavam, sempre 
achei que programação ou desenvolvimento de software fosse algo que 
só gênios conseguiam fazer.
\\\\
Aos 17 anos enquanto estava preso a um trabalho que eu 
já não me sentia mais interessado em fazer, e que sinceramente estava-
me desgastando, acabei por novamente dar de cara com o mundo da 
programação ao ver uma notícia de que Harvard estava oferecendo um curso 
completamente online e gratuito chamado 'CS50', vendo aquilo acabei por 
dar uma oportunidade a mim mesmo e experimentar.
\\\\
Por fim, depois de quase 3 anos, posso dizer que começar a 
estudar programação foi provavelmente uma das melhores escolhas que fiz 
na vida, além de claro, ter sido extremamente divertido e desafiador. 
Hoje em dia, não só tenho certeza que quero trabalhar na área mas 
também que desejo ser plenamente capaz de repassar tudo que aprendi 
para o mundo. Portanto, espero que esse e-book possa ser o meu primeiro 
grande passo para divulgar o meu conhecimento, que apesar de não muito, 
pode ajudar quem está começando.

\section{Sobre o livro}

O intuito desse e-book é ajudar iniciantes na linguagem C, que se 
sentem perdidos e não sabem por onde começar. Quando comecei a estudar sobre 
programação percebi de cara que maior parte do conteúdo bom estava disponível 
apenas em inglês, o que não foi um problema para mim que tive uma educação 
que me auxiliou e permitiu absorver esse tipo de conteúdo, porém, reconheço 
que nem todos possuíram a mesma oportunidade de aprender inglês ao decorrer 
da vida.
\\\\
Aqui você será guiado desde a como preparar a sua máquina para começar 
a programar, a até algumas dicas e truques mais avançados sobre como a linguagem 
e o compilador se comportam. Portanto, espero que esse material seja suficiente 
para preencher todas as lacunas de conhecimento necessárias para poder começar 
a programar em C sem muitos problemas.

\chapter{Introdução a linguagem C}
\clearpage

\section{O que é a linguagem C}
\subsection{História e contexto}

A linguagem C foi criada em 1970 por Dennis Ritchie com a finalidade de 
desenvolver sistemas operacionais, drivers e outros softwares de baixo 
nível. Apesar de ter mais de 60 anos, continua sendo amplamente 
utilizada em diversos tipos de hardware, desde supercomputadores até 
microcontroladores e sistemas embarcados.
\\\\
Com compiladores disponíveis em praticamente todos os sistemas modernos, 
destacou-se pela portabilidade e foi amplamente empregada na construção 
de utilitários e aplicações complexas que exigiam maior velocidade e 
eficiência. Em 1978, foi publicado o livro "The C Programming Language", 
escrito pelo coautor da linguagem C, que por muito tempo serviu como 
referência padrão para a linguagem.
\\\\
Por ser uma linguagem imperativa, com suporte à programação estruturada, 
amplo escopo lexical e acesso eficiente à memória do computador em 
baixo nível, sendo frequentemente a escolha ideal para projetos em 
sistemas com grandes restrições de hardware.

\subsection{Benefícios da linguagem e aplicações práticas}

Devido ao seu alto desempenho como linguagem compilada de baixo 
nível, a programação em C possui vantagens e desvantagens. Entre 
as vantagens, destacam-se:

\begin{itemize}
\item Alta portabilidade
\item Alto desempenho
\item Gerenciamento eficiente de memória
\item Acesso a recursos do sistema
\end{itemize}

No entanto, essas vantagens têm um preço, ideia essa que pode ser 
entendida melhor por uma frase que gosto bastante (a mesma encontrada 
em inglês no verso do livro).

\begin{quote}
    "No começo todos queremos resultados, mas no fim, tudo que queremos é controle"
    \footnote{Eskil Steenberg}
\end{quote}

Essa frase se mostra verdadeira, pois apesar de ter total controle 
e poder sobre as funcionalidades da linguagem, quando mal utilizadas 
elas podem criar vulnerabilidades de segurança e uma complexidade 
excessiva, que tende a aumentar ao longo de um projeto.
\\\\
Apesar dos desafios, quando dominada e compreendida adequadamente, 
a linguagem C não apenas oferece um entendimento sólido dos fundamentos 
da programação, mas também encontra aplicação prática no desenvolvimento de:

\begin{itemize}
\item Sistemas operacionais
\item Microcontroladores
\item Drivers
\item Linguagens de programação
\item Utilitários
\end{itemize}

\section{Ambiente de desenvolvimento}

\subsection{Configuração do ambiente}

Configurar o ambiente de desenvolvimento para programar em C é crucial para 
iniciar projetos com eficiência. Abaixo, são fornecidos passos para configurar 
o ambiente em sistemas Linux e Windows.

\subsubsection{Linux}

Para configurar o ambiente de desenvolvimento C em sistemas Linux, siga estes passos:

\begin{enumerate}
    \item \textbf{Instalação do compilador C (GCC):} 
        Em maior parte das distribuições Linux o GCC já está disponível por padrão, 
        mas caso não esteja, ele pode ser facilmente instalado utilizando o gerenciador 
        de pacotes da sua distribuição (como apt para Ubuntu/Debian ou yum para Fedora).
        \\\\
        Segue abaixo uma lista de comandos para instalar o GCC nas distribuições Linux 
        mais famosas:

        \begin{verbatim}
            # Debian
            sudo apt update
            sudo apt install build-essential
            # Red Hat-based (Fedora, CentOS)
            sudo dnf update
            sudo dnf install gcc
            # Arch Linux
            sudo pacman -Syu
            sudo pacman -S gcc
            # openSUSE
            sudo zypper refresh
            sudo zypper install gcc
            # Alpine Linux
            sudo apk update
            sudo apk add build-base

            # Para checar se o GCC foi instalado 
            # corretamente use o comando abaixo
            gcc --version
        \end{verbatim}
\end{enumerate}

\subsubsection{Windows}

Para configurar o ambiente de desenvolvimento C em sistemas Windows, siga 
estes passos:

\begin{enumerate}
    \item \textbf{Instalação do compilador C (GCC):} 
        No Windows a utilização de ferramentas como o GCC acaba por ser 
        um pouco mais complexa pois é necessária a instalação do MinGW ou 
        MSYS2. Essas ferramentas fornecem um ambiente similar ao Linux 
        para compilar código C.

        Para instalar o MinGW, siga os passos abaixo:

        \begin{enumerate}
            \item Acesse o site oficial do MinGW: \texttt{https://osdn.net/projects/mingw/}
            \item Baixe o instalador adequado para o seu sistema (32 ou 64 bits).
            \item Execute o instalador baixado (\texttt{mingw-w64-install.exe}).
            \item Durante a instalação, selecione os pacotes desejados. O 
                pacote essencial para o GCC é o \texttt{mingw32-gcc} para 
                sistemas 32 bits ou \texttt{mingw64-gcc} para sistemas 64 bits.
            \item Siga as instruções do instalador para concluir a instalação.
        \end{enumerate}

        Após a instalação, é necessário configurar o ambiente para usar o GCC via linha de comando:

        \begin{enumerate}
            \item Adicione o diretório bin do MinGW ao caminho do sistema:
                \begin{itemize}
                    \item Para Windows 10, pesquise "Editar variáveis de ambiente 
                        do sistema" na barra de pesquisa e clique em "Variáveis 
                        de ambiente". Em "Variáveis do sistema", selecione a 
                        variável PATH e clique em "Editar...". Adicione o caminho 
                        para o diretório bin do MinGW (exemplo: 
                        \texttt{C:\textbackslash MinGW\textbackslash bin}) e 
                        clique em OK.
                    \item Para versões anteriores do Windows, o processo é 
                        similar, mas pode variar um pouco, caso você utilize 
                        uma versão mais antiga do sistema talvez seja necessário 
                        que você faça uma pesquisa a parte para finalizar o 
                        processo corretamente.
                \end{itemize}
            \item Verifique se a instalação foi bem-sucedida abrindo o Prompt de Comando (cmd) e digitando \texttt{gcc -v}. Deverá mostrar a versão do GCC instalada.
        \end{enumerate}
    
    \item \textbf{Editor de texto ou IDE:} 
        Escolha um editor de texto ou uma IDE (Ambiente de Desenvolvimento Integrado), 
        para escrever e compilar código C. Alguns exemplos populares são o Visual 
        Studio Code, Atom, Sublime Text, ou IDEs como o Code::Blocks ou o Eclipse 
        com plugins para C/C++.

        {\tiny Eu pessoalmente recomendo a utilização do Nvim que apesar de ser 
        uma ferramenta de difícil utilização, pode acabar por aumentar muito a sua 
        produtividade no médio/longo prazo.}
\end{enumerate}

Se você seguiu os passos acima com sucesso, seu ambiente de desenvolvimento 
C está configurado e pronto para uso. Agora é o momento de escrever seu 
primeiro programa!

\subsection{Primeiro programa}

\section{Conceitos Fundamentais}
\subsection{Operadores}
\subsection{Variáveis e tipos de dados}
\subsection{Conversão de tipos}
\subsection{Loops e condicionais}
\subsection{Funções e bibliotecas}
\subsection{Escopo}
\subsection{Entrada e saída de dados}

\chapter{Estrutura de dados e algoritmos em C}

\section{Ponteiros}
\subsection{Como a memória funciona}
\subsection{Endereço de memória}
\subsection{Alocação de memória}
\subsection{Heap}
\subsection{Stack}

\section{Strings}
\subsection{Por baixo dos panos}
\subsection{Declaração}
\subsection{Manipulação}

\chapter{Tópicos avançados}

\section{Manipulando arquivos}
\subsection{Leitura e escrita de arquivos}

\section{Manipulação de bits}
\subsection{NOT}
\subsection{AND}
\subsection{OR}
\subsection{XOR}

\chapter{Depuração e testes}

\section{Depuração}
\subsection{Estratégias de depuração de código}
\subsection{GDB}
\subsection{Valgrind}

\section{Testes}
\subsection{Implementando testes}

\chapter{Considerações finais}
\subsection{Dicas finais}
\subsection{Recursos e referências adicionais}
\subsection{Mudanças no padrão C2X}
\subsection{Obrigado}

\end{document}
